\subsubsection*{Propositional Inference Rules}

\nbgpara{5}{Conjunction}

\nbgparar{$\bm{\land}$-I} \quad Conjunction is governed by the following introduction rule:

\begin{prooftree}
    \AxiomC{\(\varphi\)}
    \AxiomC{\(\psi\)}
    \RightLabel{\(\land\)-I}
    \BinaryInfC{\(\varphi\land\psi\)}
\end{prooftree}

\vspace{0.3cm}

\nbgparar{$\bm{\land}$-E} \quad Conjunction is governed by the following two elimination rules:

\begin{center}
\begin{minipage}{0.55\textwidth}
\centering
\begin{prooftree}
    \AxiomC{\(\varphi\land\psi\)}
    \RightLabel{\(\land\)-E}
    \UnaryInfC{\(\varphi\)}
\end{prooftree}
\end{minipage}
\hfill
\begin{minipage}{0.42\textwidth}
\centering
\begin{prooftree}
    \AxiomC{\(\varphi\land\psi\)}
    \RightLabel{\(\land\)-E}
    \UnaryInfC{\(\psi\)}
\end{prooftree}
\end{minipage}
\end{center}

\vspace{0.5cm}

\nbgpara{6}{Disjunction}

\nbgparar{$\bm{\lor}$-I} \quad Disjunction is governed by the following two introduction rules:

\begin{center}
\begin{minipage}{0.55\textwidth}
\centering
\begin{prooftree}
    \AxiomC{\(\varphi\)}
    \RightLabel{\(\lor\)-I}
    \UnaryInfC{\(\varphi\lor\psi\)}
\end{prooftree}
\end{minipage}
\hfill
\begin{minipage}{0.42\textwidth}
\centering
\begin{prooftree}
    \AxiomC{\(\psi\)}
    \RightLabel{\(\lor\)-I}
    \UnaryInfC{\(\varphi\lor\psi\)}
\end{prooftree}
\end{minipage}
\end{center}

\vspace{0.3cm}

\nbgparar{$\bm{\lor}$-E} \quad Disjunction is governed by the following elimination rule:

\begin{prooftree}
    \AxiomC{\(\varphi\lor\psi\)}

    \AxiomC{}
    \RightLabel{\(n\)}
    \UnaryInfC{\(\varphi\)}
    \noLine
    \UnaryInfC{\(\vdots\)}
    \noLine
    \UnaryInfC{\(\chi\)}

    \AxiomC{}
    \RightLabel{\(m\)}
    \UnaryInfC{\(\psi\)}
    \noLine
    \UnaryInfC{\(\vdots\)}
    \noLine
    \UnaryInfC{\(\chi\)}

    \RightLabel{\(\lor\)-E, \(n,m\)}
    \TrinaryInfC{\(\chi\)}
\end{prooftree} \vspace{0.3cm}

The labels \(n\) and \(m\) are distinct. The application discharges the assumption occurrences marked with \(n\) only in the left minor premise subderivation and the assumption occurrences marked with \(m\) only in the right minor premise subderivation.

\vspace{0.3cm}

\nbgpara{7}{Implication}

\nbgparar{$\bm{\to}$-I} \quad Implication is governed by the following introduction rule:

\begin{prooftree}
    \AxiomC{}
    \RightLabel{\(n\)}
    \UnaryInfC{\(\varphi\)}
    \noLine
    \UnaryInfC{\(\vdots\)}
    \noLine
    \UnaryInfC{\(\psi\)}
    \RightLabel{\(\to\)-I, \(n\)}
    \UnaryInfC{\(\varphi\to\psi\)}
\end{prooftree}

The application discharges every assumption occurrence marked with \(n\) in its premise subderivation.

An application of \(\to\)-I may be vacuous: the rule may be applied even if its premise subderivation contains no assumption occurrence marked with the indicated label \(n\):

\begin{prooftree}
    \AxiomC{}
    \noLine
    \UnaryInfC{\(\vdots\)}
    \noLine
    \UnaryInfC{\(\psi\)}
    \RightLabel{\(\to\)-I, \(n\)}
    \UnaryInfC{\(\varphi\to\psi\)}
\end{prooftree}

\vspace{0.3cm}

\nbgparar{$\bm{\to}$-E} \quad Implication is governed by the following elimination rule:

\begin{prooftree}
    \AxiomC{\(\varphi\to\psi\)}
    \AxiomC{\(\varphi\)}
    \RightLabel{\(\to\)-E}
    \BinaryInfC{\(\psi\)}
\end{prooftree}

\nbgpara{8}{Biconditional}

\nbgparar{$\bm{\leftrightarrow}$-I} \quad The biconditional is governed by the following introduction rule:

\begin{prooftree}
    \AxiomC{\(\varphi\to\psi\)}
    \AxiomC{\(\psi\to\varphi\)}
    \RightLabel{\(\leftrightarrow\)-I}
    \BinaryInfC{\(\varphi\leftrightarrow\psi\)}
\end{prooftree}

\vspace{0.3cm}

\nbgparar{$\bm{\leftrightarrow}$-E} \quad The biconditional is governed by the following two elimination rules:

\begin{center}
\begin{minipage}{0.55\textwidth}
\centering
\begin{prooftree}
    \AxiomC{\(\varphi\leftrightarrow\psi\)}
    \AxiomC{\(\varphi\)}
    \RightLabel{\(\leftrightarrow\)-E}
    \BinaryInfC{\(\psi\)}
\end{prooftree}
\end{minipage}
\hfill
\begin{minipage}{0.42\textwidth}
\centering
\begin{prooftree}
    \AxiomC{\(\varphi\leftrightarrow\psi\)}
    \AxiomC{\(\psi\)}
    \RightLabel{\(\leftrightarrow\)-E}
    \BinaryInfC{\(\varphi\)}
\end{prooftree}
\end{minipage}
\end{center} 

\nbgpara{9}{Negation}

\nbgparar{$\bm{\neg}$-I} \quad Negation is governed by the following introduction rule:

\begin{prooftree}
    \AxiomC{}
    \RightLabel{\(n\)}
    \UnaryInfC{\(\varphi\)}
    \noLine
    \UnaryInfC{\(\vdots\)}
    \noLine
    \UnaryInfC{\(\bot\)}
    \RightLabel{\(\neg\)-I, \(n\)}
    \UnaryInfC{\(\neg\varphi\)}
\end{prooftree}

\vspace{0.3cm}

\nbgparar{$\bm{\neg}$-E} \quad Negation is governed by the following elimination rule:

\begin{prooftree}
    \AxiomC{\(\varphi\)}
    \AxiomC{\(\neg\varphi\)}
    \RightLabel{\(\neg\)-E}
    \BinaryInfC{\(\bot\)}
\end{prooftree}


\nbgpara{10}{Falsity}

\nbgparar{$\bm{\bot}$-E} \quad From falsity, any formula may be inferred:

\begin{prooftree}
    \AxiomC{\(\bot\)}
    \RightLabel{\(\bot\)-E}
    \UnaryInfC{\(\varphi\)}
\end{prooftree}

\vspace{0.3cm}

\nbgpara{11}{Classical Reductio}

\nbgparar{CR} \quad Classical reductio is governed by the following rule:

\begin{prooftree}
    \AxiomC{}
    \RightLabel{\(n\)}
    \UnaryInfC{\(\neg\varphi\)}
    \noLine
    \UnaryInfC{\(\vdots\)}
    \noLine
    \UnaryInfC{\(\bot\)}
    \RightLabel{CR, \(n\)}
    \UnaryInfC{\(\varphi\)}
\end{prooftree} \vspace{0.3cm}

The propositional rules stated in paragraphs 5--10 are intuitionistically valid. The addition of CR is what makes the propositional fragment of the deductive system classical. \vspace{-0.3cm}

